\phase{7 --- Apple execution and algorithm implementations}

\textbf{Objective.} Deliver a production-quality Apple path without changing the public image-processing contract.

\deliverables
\begin{enumerate}
  \item CPU acceleration adapters using Accelerate/vImage where they satisfy the operation contract.
  \item Core Image adapters for appropriate image operations, with explicit handling of its context, color management, deferred evaluation, and materialization boundaries.
  \item A native Metal provider implementing \GEP{} on supported Apple devices.
  \item Metal operation implementations for the initial operation set where a native kernel offers a measured advantage or an essential capability.
  \item Safe interoperation among CPU memory, Core Image representations, and Metal resources where documented and supported; otherwise use explicit copy/materialization adapters.
  \item Headless command-line build/test scripts. The workflow must not require an Xcode project or the Xcode IDE; native Metal compilation still requires a compatible Apple SDK/toolchain and must be validated on the actual target environment.
\end{enumerate}

\exitcriteria
\begin{itemize}
  \item CPU and Apple implementations pass the same operation conformance tests.
  \item Resource lifetimes and synchronization are verified under repeated asynchronous execution.
  \item No unintended CPU/GPU round trip is present in a same-device sequence, as demonstrated by instrumentation.
  \item Core Image's deferred results are not reported as materialized outputs before the public completion contract is met.
  \item The public operation and image APIs do not expose Core Image or Metal types.
\end{itemize}

\phase{8 --- OpenCV and C/C++ interoperability adapters}

\textbf{Objective.} Reuse OpenCV's mature algorithms without making OpenCV the library's semantic or storage model.

\deliverables
\begin{enumerate}
  \item A narrow interop boundary for calling C++ implementations from the public library (a C ABI or an appropriate platform bridge, chosen after validating supported toolchains).
  \item Explicit adapters for wrapping/copying image memory, including lifetime, stride, alignment, ownership, and mutability.
  \item OpenCV CPU operation implementations registered as \IPP{} candidates only when their semantics match the operation contract.
  \item Optional OpenCV CUDA/OpenCL adapters treated as algorithm implementations, not as replacements for the general provider protocols.
  \item Diagnostics that disclose when an adapter copies, converts, or changes residency.
\end{enumerate}

\exitcriteria
\begin{itemize}
  \item No C++ type such as \texttt{cv::Mat}, \texttt{cv::UMat}, or \texttt{cv::cuda::GpuMat} leaks into the public semantic API.
  \item External buffer lifetime and ownership are tested, including asynchronous use.
  \item OpenCV implementations pass the same operation conformance suite; unsupported semantic modes are rejected.
  \item Tests demonstrate whether the chosen adapter path shares storage or performs a copy.
\end{itemize}

\phase{9 --- CUDA and OpenCL native providers}

\textbf{Objective.} Add direct implementations of the general execution contract for NVIDIA CUDA and OpenCL, then register operation kernels independently.

\deliverables
\begin{enumerate}
  \item CUDA provider: device discovery, contexts/streams as applicable, resource allocation/import, submission, events/completion, and safe release.
  \item OpenCL provider: platform/device discovery, contexts, queues, memory objects, event dependencies, and safe release.
  \item Native operation implementations for a subset selected by benchmark evidence.
  \item Interop contracts for OpenCV CUDA/OpenCL adapters only where real sharing and synchronization guarantees exist.
  \item CI/test plans using actual supported NVIDIA and OpenCL hardware; mock-only qualification is insufficient.
\end{enumerate}

\exitcriteria
\begin{itemize}
  \item Both providers pass the generic \GEP{} conformance suite.
  \item Every registered operation implementation passes the \IPP{} operation suite.
  \item Provider switching preserves output semantics and reports transfers/conversions accurately.
  \item Device loss, allocation failure, cancellation, and synchronization errors do not corrupt image state.
  \item CPU fallback occurs only when it is qualified and allowed by request constraints.
\end{itemize}
